% ===== 第一段：研究背景引言 =====
\section{立题依据及价值}

\subsection{研究背景与有关研究现状}

近年来，大语言模型（Large Language Model，LLM）逐步从"文本生成模型"演化为能够进行任务规划、工具调用、环境交互和多轮反馈修正的智能体（Agent）。
在ReAct\cite{yao2023react}、Toolformer\cite{schick2023toolformer}、ToolLLM\cite{qin2024toolllm}、Gorilla\cite{patil2024gorilla}等代表性工作中，
模型不再仅依赖内部参数完成推理，而是通过外部API、检索系统、文件系统、邮件日程、数据库、浏览器或代码执行器等工具获取信息并改变外部环境。
CodeAct\cite{wang2024codeact}则进一步表明，将工具调用以可执行代码动作的形式组织，可以显著提升Agent在多步任务上的完成能力。
Wang等\cite{wang2024survey_agents}和Xi等\cite{xi2025rise}对LLM Agent的能力框架、应用场景和挑战进行了系统综述，
揭示Agent系统已逐步进入办公自动化、软件工程、网络安全运维、数据分析和智能物联网等关键应用领域\cite{park2023generativeagents}；
Weidinger等\cite{weidinger2022taxonomy}则从更宏观的视角对语言模型的风险进行了系统分类，指出随着模型与外部环境交互能力的增强，安全与治理挑战将成倍增加。

然而，工具调用能力的引入从根本上改变了LLM的安全威胁面。
Chu等\cite{chu2026systematic}指出，与传统LLM的文本生成风险不同，工具调用的每一次执行都构成具有现实侧效应的动作，
且工具返回结果通常在无显式信任标记的情况下被注入Agent规划上下文，使得数据流污染、权限误用和不可逆操作等问题从推理层延伸至执行层；
尤其是"语义权限提升"问题——每次工具调用均经过授权，但授权动作的语义组合可能超越预期授权范围，而这一过程对基于单次调用的访问控制系统完全不可见。
Deng等\cite{deng2025agentsunderthreat}的综述则进一步表明，序列规划中的单次错误会沿推理链持续放大与传播，最终导致级联失效——
这意味着Agent工具调用的安全性不能通过孤立评估每次调用来保证，而必须从调用链的整体结构出发建立风险度量体系。

从工具调用攻防现状看，WebArena\cite{zhou2023webarena}、Mind2Web\cite{deng2023mind2web}、OSWorld\cite{xie2024osworld}、
Windows Agent Arena\cite{bonatti2024windowsarena}等基准显示，真实环境中的Agent任务通常呈现长时序、跨应用和多工具组合特征。
Dziemian等\cite{dziemian2026vulnerable}在迄今最大规模的红队竞赛（覆盖22个前沿Agent、44个真实部署场景）中发现，
当前Agent系统在多种模型和部署场景下攻击成功率接近100\%，且成功攻击具有高度可迁移性，而Agent鲁棒性与模型规模、能力或推理算力之间几乎不存在相关性——
说明单纯提升模型能力无法解决工具调用链的安全问题，需要独立的审计与度量机制。
AgentDojo\cite{balunovic2024agentdojo}作为动态评测框架，系统量化了提示注入对工具调用型Agent的威胁；
Imprompter\cite{fu2024imprompter}、CaMeL\cite{debenedetti2025defeating}、FIDES\cite{costa2025fides}、
Firewalled Agentic Networks\cite{abdelnabi2025firewalls}等工作进一步表明，
工具调用过程中的提示注入、数据流污染、权限误用和轨迹偏离已成为Agent安全的核心问题。
然而，上述工作的评价视角大多停留于"攻击是否成功"或"防御是否阻断"，
对工具调用链本身的风险结构、过程级可解释评估和事后归因机制关注不足。

从多智能体安全和系统安全的视角看，将工具调用执行轨迹结构化为溯源图（provenance graph），是解决上述问题的天然路径。
传统系统安全领域长期使用审计日志和溯源图刻画进程、文件、网络连接之间的因果关系\cite{gao2021threatraptor,cheng2024kairos,inam2023sok}；
对于LLM Agent，prompt、plan、tool call、observation、memory、final answer等执行片段同样构成可追踪的行为链条。
将这些片段图化，可以同时表达"谁触发了哪个工具""参数来源于何处""敏感数据经过哪些转换""最终风险由哪些路径贡献"等信息，
从而为工具滥用检测与风险归因提供结构化基础\cite{wang2026traces,dewitt2026openchallenges}。

本课题拟围绕"工具调用链风险度量"和"行为轨迹图检测归因"两个研究点展开，
从序列层面建立可解释风险指标体系，从图结构层面实现工具滥用检测与责任归因，
以下分四个方向梳理国内外相关研究现状。
